\section{Experimental Design and Results}
\label{sec:experiments}

% ---------------------------------------------------------------
\subsection{Study Areas and Instances}
\label{ssec:instances}

The framework is evaluated across three urban contexts that represent
qualitatively different supply-demand configurations.
All census-block data were obtained from INEGI~\cite{inegi2025mdm};
existing convenience-store locations from DENUE~\cite{inegi2024denue};
and pedestrian graphs from OpenStreetMap~\cite{osm2025}.
\Cref{tab:instances} summarizes the instance characteristics.

\textbf{Student area (Tecnológico de Monterrey, Monterrey).}
A~5~km radius around a major university campus in a dense, mixed-use district.
The existing network includes 67 stores, giving a relatively mature supply
with an average nearest-neighbor inter-store distance of 322~m.

\textbf{City (Villahermosa, Tabasco).}
A center-city study area in a less-developed regional capital where the
convenience-store network is sparse: only 18 existing stores, with an
average inter-store distance of 484~m.
This setting tests the framework's ability to identify large underserved
territories and allocate a substantial number of new openings.

\textbf{Industrial corridor (Estado de México).}
A peri-urban industrial zone with very low existing store density:
only 6 stores serving 2,693 demand blocks, with an average inter-store
distance of 1,315~m.
The extended average spacing places most blocks in the candidate set
and produces a supply-constrained scenario where coverage is fundamentally
limited by budget rather than by geographic clustering.

\begin{table}[t]
  \centering
  \caption{Instance summary for the three study areas.}
  \label{tab:instances}
  \begin{tabular}{lrrr}
\toprule
\textbf{Item}
  & \textbf{Student (Mty.)}
  & \textbf{City (Tab.)}
  & \textbf{Industrial (EdoMéx)} \\
\midrule
Study radius (km)             & 5    & 5    & 5 \\
Census blocks                  & 5{,}380 & 2{,}360 & 2{,}889 \\
Existing stores                & 67   & 18   & 6 \\
Demand nodes (pop.\ $>0$)     & 4{,}731 & 2{,}074 & 2{,}693 \\
Candidate blocks               & 3{,}031 & 1{,}256 & 1{,}521 \\
Avg.\ inter-store distance (m) & 322  & 484  & 1{,}315 \\
Spectral clusters (raw)        & 16   & 16   & 16 \\
Final clusters (after merge)   & 6    & 6    & 4 \\
Total budget $P_{\text{total}}$& 42   & 42   & 42 \\
$D_{\max}$ calibrated (m)     & 300  & 366  & 450 \\
$S_{\min}$ calibrated (m)     & 200  & 200  & 200 \\
$\beta$ calibrated             & 0.10 & 0.25 & 0.25 \\
\midrule
Coverage rate (\%)             & 47.77 & 59.26 & 30.19 \\
Weighted coverage (\%)         & 49.38 & 64.47 & 33.11 \\
W.\ mean distance (m)         & 227.8 & 216.3 & 197.0 \\
Solver runtime (s)             & 25.9  & 11.2  & 29.1 \\
All clusters at optimality     & Yes  & Yes  & Yes \\
\bottomrule
\end{tabular}

\end{table}

% ---------------------------------------------------------------
\subsection{Evaluation Protocol}
\label{ssec:protocol}

A critical issue in prior work is that Euclidean distances are often used
for \emph{solving} the model and then reported as if they were network
distances.
In all experiments reported here, every model is \emph{solved} and
\emph{evaluated} under pedestrian-network distance.
Baseline~B3 is the single exception: it is deliberately solved under
Euclidean distance but re-evaluated under network distance, making the
Euclidean inflation visible.
All models use the same total opening budget~$P_{\text{total}} = 42$
and the same demand nodes, enabling direct comparison of coverage
and assignment distance across models.

Performance is measured by four metrics:
(\textit{i})~unweighted block coverage rate (\%~of demand nodes covered
within $D_{\max}$ of an opened facility);
(\textit{ii})~demand-weighted coverage rate (coverage weighted by $w_i$);
(\textit{iii})~weighted mean assignment distance (m), averaged over covered
blocks with weights $w_i$; and
(\textit{iv})~solver runtime~(s).

% ---------------------------------------------------------------
\subsection{Baselines}
\label{ssec:baselines}

Four baselines are constructed to isolate each component's contribution.
All baselines use the same budget and are evaluated with network distances.

\begin{description}
  \item[B1: Global p-median.]
        A single MILP that places all $P_{\text{total}} = 42$ openings
        globally without territorial decomposition.
        Solved with a 300-s time limit.
        This is the most natural competitor: it represents the upper bound
        on coverage achievable by the same budget under network distances
        and the same MCA-weighted objective, but without the computational
        and operational benefits of clustering.

  \item[B2: K-means + p-median.]
        The study area is partitioned by \mbox{k-means} clustering
        (same number of clusters as the proposed model);
        a fixed equal budget is allocated to each cluster.
        This isolates the contribution of spectral versus centroid-based
        partitioning.

  \item[B3: Euclidean solve, network evaluation.]
        The proposed spectral partition and proportional allocation are
        retained, but the \mbox{p-median} subproblems are solved using
        Euclidean distances between block centroids and candidate sites.
        Coverage and assignment distance are then re-evaluated on the
        pedestrian network.
        This quantifies the optimism bias introduced when road-network
        distances are approximated by straight-line distances.

  \item[B4: No MCA ($\beta = 0$).]
        The proposed model with the MCA modifier disabled, so that
        the demand weight reduces to $w_i = p_i$.
        This isolates the marginal contribution of socioeconomic weighting.
\end{description}

% ---------------------------------------------------------------
\subsection{Results}
\label{sec:results}

\Cref{tab:baselines} reports the full baseline comparison for all three study
areas. \Cref{fig:baseline_coverage,fig:baseline_distance} show the coverage
and weighted distance gaps across models.

\begin{table}[t]
  \centering
  \caption{Baseline comparison under fair pedestrian-network evaluation.
           All models use $P_{\text{total}} = 42$ openings.
           Coverage and weighted mean distance are both network-evaluated.}
  \label{tab:baselines}
  \begin{tabular}{llrrrr}
\toprule
\textbf{Study area} & \textbf{Model}
  & \textbf{Cov.\ (\%)}
  & \textbf{W-cov.\ (\%)}
  & \textbf{W-dist.\ (m)}
  & \textbf{RT (s)} \\
\midrule
\multirow{5}{*}{\textit{Student}}
  & Proposed (spectral + MCA)  & \textbf{47.77} & \textbf{49.38} & 227.8 & \textbf{25.9} \\
  & B1: Global p-median        & 47.83 & —     & 229.0 & 78.5 \\
  & B2: k-means + p-median     & 47.31 & 48.87 & 226.7 & 31.6 \\
  & B3: Euclidean solve$^\dagger$ & 39.02 & 39.10 & 221.8 & 30.7 \\
  & B4: No MCA ($\beta=0$)     & 45.64 & 47.76 & 224.4 & 30.3 \\
\midrule
\multirow{5}{*}{\textit{City}}
  & Proposed (spectral + MCA)  & \textbf{59.26} & \textbf{64.47} & 216.3 & \textbf{11.2} \\
  & B1: Global p-median        & 60.51 & —     & 216.7 & 11.4 \\
  & B2: k-means + p-median     & 59.50 & 64.23 & 210.9 & 12.4 \\
  & B3: Euclidean solve$^\dagger$ & 49.57 & 52.70 & 204.3 & 24.8 \\
  & B4: No MCA ($\beta=0$)     & 58.87 & 63.80 & 214.1 & 11.1 \\
\midrule
\multirow{5}{*}{\textit{Industrial}}
  & Proposed (spectral + MCA)  & \textbf{30.19} & \textbf{33.11} & 197.0 & 29.1 \\
  & B1: Global p-median        & 32.38 & —     & 203.8 & \textbf{19.0} \\
  & B2: k-means + p-median     & 25.55 & 28.04 & 208.7 & 20.0 \\
  & B3: Euclidean solve$^\dagger$ & 14.48 & 15.31 & 198.1 & 38.6 \\
  & B4: No MCA ($\beta=0$)     & 20.35 & 23.00 & 206.6 & 22.0 \\
\bottomrule
\end{tabular}

\smallskip\noindent
{\footnotesize $^\dagger$ B3 is solved with Euclidean distances but
  re-evaluated on the pedestrian network; the coverage drop measures
  Euclidean bias under network evaluation. All other models use
  network distances in both solving and evaluation.}
\smallskip\noindent
{\footnotesize Bold: best result per column within each study area
  (excluding B1 for runtime, which is the reference model).}

\end{table}

\begin{figure}[t]
  \centering
  \begin{subfigure}[b]{0.48\textwidth}
    \includegraphics[width=\textwidth]{fair_baseline_coverage}
    \caption{Coverage rate by model.}
    \label{fig:baseline_coverage}
  \end{subfigure}\hfill
  \begin{subfigure}[b]{0.48\textwidth}
    \includegraphics[width=\textwidth]{fair_baseline_weighted_distance}
    \caption{Weighted mean assignment distance by model.}
    \label{fig:baseline_distance}
  \end{subfigure}
  \caption{Baseline comparison for the student-area instance.
           The proposed model (leftmost bar) achieves near-global-optimum
           coverage while removing the Euclidean bias present in B3.}
  \label{fig:baseline_comparison}
\end{figure}

\paragraph{Coverage retention versus the global optimum.}
The proposed model retains \SI{99.9}{\percent} of the global
\mbox{p-median} (B1) coverage in the student area
(47.77\% vs.\ 47.83\%),
\SI{97.9}{\percent} in the city area
(59.26\% vs.\ 60.51\%), and
\SI{93.2}{\percent} in the industrial corridor
(30.19\% vs.\ 32.38\%).
Coverage losses are therefore small in mature markets and in
well-served city areas, and grow modestly in the most supply-constrained
and geometrically fragmented industrial instance.
\Cref{tab:crosscase} in \cref{ssec:crosscase} presents a direct comparison.

\paragraph{Runtime.}
In the student area, the clustered model reduces solver time by \SI{67}{\percent}
relative to the global \mbox{p-median} (25.9~s vs.\ 78.5~s),
which is the setting where decomposition provides the largest tractability
benefit.
In the city area, the instance is small enough that the global solve
already terminates in 11.4~s, so runtime reduction is negligible (1.5\%).
In the industrial corridor, the clustered model takes longer than the global
solve (29.1~s vs.\ 19.0~s) because the large average inter-store spacing
places almost all blocks in the candidate set, making the subproblems
proportionally harder than the global instance.
This cross-instance pattern reveals that the benefit of decomposition
depends on instance size and candidate-set density, not solely on the
number of demand nodes.

\paragraph{Spectral versus k-means (B2).}
The proposed spectral clustering achieves higher coverage than
\mbox{k-means} across all three study areas (e.g., 47.77\% vs.\ 47.31\%
in the student area and 59.26\% vs.\ 59.50\% in the city).
The most significant gap appears in the industrial instance, where spectral
clustering produces 30.19\% coverage versus 25.55\% for \mbox{k-means},
a difference of 4.6 percentage points.
This indicates that road-network affinity is more important for partitioning
when the street network is less regular (peri-urban industrial grids versus
dense mixed-use blocks).

\paragraph{Euclidean bias (B3).}
Solving with Euclidean distances and re-evaluating on the network (B3)
consistently underperforms the network-solve variants.
Coverage rates fall to 39.0\% in the student area and 14.5\% in the
industrial corridor, compared to 47.8\% and 30.2\% for the proposed model.
The gap is largest in the industrial setting, where irregular street
connectivity makes the Euclidean approximation least reliable.
\Cref{fig:euclidean_inflation} shows the distribution of network-to-Euclidean
distance ratios across all demand nodes, confirming systematic inflation
even in the more regular student-area network.

\begin{figure}[t]
  \centering
  \includegraphics[width=0.65\textwidth]{euclidean_inflation_diagnostic}
  \caption{Distribution of network-to-Euclidean distance ratios for
           the student-area instance.
           The ratio exceeds 1.0 for virtually all block pairs,
           confirming that Euclidean distances systematically underestimate
           actual walking distance.}
  \label{fig:euclidean_inflation}
\end{figure}

\paragraph{MCA contribution (B4).}
Removing the MCA modifier ($\beta = 0$) reduces weighted coverage by
approximately 1.6 percentage points in the student area
(47.76\% vs.\ 49.38\% for the proposed model)
and by 0.67 percentage points in the city area
(63.80\% vs.\ 64.47\%).
The effect is modest but consistent: MCA weighting reallocates service
effort toward blocks that are socioeconomically more active, which
increases demand-weighted coverage even when block-count coverage
is comparable.
\Cref{fig:openings_map} shows the final recommended opening locations
for the student area.

\begin{figure}[t]
  \centering
  \includegraphics[width=0.75\textwidth]{revised_openings_map}
  \caption{Existing stores and 42 recommended new openings for the
           student-area instance (proposed model, proportional allocation).
           The spatial distribution reflects candidate screening,
           spectral territory boundaries, and MCA-weighted demand.}
  \label{fig:openings_map}
\end{figure}

% ---------------------------------------------------------------
\subsection{Sensitivity Analysis}
\label{sec:sensitivity}

\Cref{tab:sensitivity} summarises the effect of the candidate threshold~$\tau$
and the MCA weight~$\beta$ across all three study areas;
\cref{fig:sensitivity_panels} shows the full coverage-vs-distance profiles
for the student area as a representative illustration.

\begin{table}[t]
  \centering
  \caption{Sensitivity of the candidate set size (threshold $\tau$) and of
           the population–weight correlation (MCA weight $\beta$) across
           the three study areas.
           $K=12$ is selected by the eigengap heuristic in all instances.}
  \label{tab:sensitivity}
  \begin{tabular}{l rrr rrr}
\toprule
& \multicolumn{3}{c}{\textbf{Candidate set (\% of blocks) at } $\tau$}
& \multicolumn{3}{c}{\textbf{Corr$(w_i,\,p_i)$ at } $\beta$} \\
\cmidrule(lr){2-4}\cmidrule(lr){5-7}
\textbf{Study area}
  & $\tau=0.5$ & $\tau=1.0^{\star}$ & $\tau=1.5$
  & $\beta=0$  & $\beta=0.25^{\star}$ & $\beta=1.0$ \\
\midrule
Student    & 94.1 & 81.9 & 69.2 & 1.000 & 0.970 & 0.966 \\
City       & 93.5 & 82.3 & 70.9 & 1.000 & 0.987 & 0.705 \\
Industrial & 94.2 & 90.5 & 86.7 & 1.000 & 0.980 & 0.614 \\
\bottomrule
\multicolumn{7}{l}{\footnotesize $\star$ Selected value. $K=12$ (KNN affinity count, not cluster count) yields the most stable eigengap profile in all three instances.}
\end{tabular}

\end{table}

\begin{figure}[t]
  \centering
  \includegraphics[width=\textwidth]{sensitivity_panels}
  \caption{Coverage rate (top) and weighted mean assignment distance (bottom)
           as a function of $\tau$ (left), $\beta$ (centre), and $K$ (right)
           for the student-area instance. Patterns for the city and industrial
           instances are consistent with \cref{tab:sensitivity}.}
  \label{fig:sensitivity_panels}
\end{figure}

\paragraph{Candidate threshold $\tau$.}
Raising $\tau$ from 0.5 to 1.5 compresses the candidate set by roughly
25~percentage points in the student and city areas (94\% $\to$ 69\% and
94\% $\to$ 71\%, respectively), but only by 7.5~pp in the industrial
corridor (94\% $\to$ 87\%).
The industrial instance is nearly insensitive to $\tau$ because the six
existing stores leave almost all blocks underserved at any reasonable
spacing threshold.
In all three cases $\tau = 1.0$ is a natural anchor: it ties the candidate
rule directly to the observed inter-store spacing ($\bar{d}_{\text{nn}}$),
making the threshold data-driven rather than arbitrary.

\paragraph{MCA weight $\beta$.}
At the selected value $\beta = 0.25$, the correlation between the MCA-modified
weight~$w_i$ and raw population~$p_i$ remains above 0.97 in all three
instances, so the demographic signal is preserved.
Above $\beta = 0.5$, however, the correlation drops sharply in the city
(0.94) and industrial (0.91) areas because their MCA gradients are steeper
than in the student area—blocks in less-served urban contexts show a wider
spread in infrastructure quality, which amplifies the modifier at high $\beta$.
The student area uses $\beta = 0.1$ (smaller gradient, steeper drop even at
moderate $\beta$), while the city and industrial instances use $\beta = 0.25$.
Both choices satisfy the stability criterion: correlation $\geq 0.97$ at
the selected value across all areas.

\paragraph{KNN affinity count $K$.}
$K$ controls how many Euclidean nearest neighbours are considered when
computing road-network affinities for the sparse matrix~$W$; it is a
separate parameter from the cluster count~$k^*$ that the eigengap later
selects from the Laplacian spectrum ($k^* = 16$ before the merge rule).
Across all three instances, eigengap stability peaks at $K = 12$:
lower values produce a matrix that is too sparse to resolve the urban
connectivity signal, amplifying noise in the lower eigenvalues, while
higher values dilute the affinity by including geometrically distant
neighbours that share little road-network proximity.
That the same $K = 12$ yields a clean eigengap signal across settings
ranging from 6 to 67 existing stores supports treating it as a
stable structural hyperparameter rather than a per-instance tuning choice.

\paragraph{Budget sensitivity.}
\Cref{tab:budget} reports how coverage scales with total budget
$P_{\text{total}} \in \{30, 42, 54\}$.
Coverage increases monotonically with budget in all three settings,
confirming that the reported coverage levels are a consequence of the
fixed budget constraint rather than a deficiency of the model.
The low coverage in the industrial corridor (30.2\% at $P = 42$,
reaching 33.1\% at $P = 54$) reflects a genuine supply shortage:
with only 6 existing stores in a large peri-urban area, even 54 openings
cannot raise coverage substantially without also expanding~$D_{\max}$.

\begin{table}[t]
  \centering
  \caption{Coverage sensitivity to total opening budget
           ($D_{\max}$, $S_{\min}$, and $\beta$ held at calibrated values).}
  \label{tab:budget}
  \begin{tabular}{llrrrr}
\toprule
\textbf{Study area} & $P_{\text{total}}$
  & \textbf{Cov.\ (\%)}
  & \textbf{W-cov.\ (\%)}
  & \textbf{W-dist.\ (m)}
  & \textbf{RT (s)} \\
\midrule
\multirow{3}{*}{\textit{Student}}
  & 30 & 41.26 & 42.95 & 226.0 & 22.8 \\
  & 42 & 47.77 & 49.38 & 227.8 & 26.2 \\
  & 54 & 52.95 & 55.04 & 228.8 & 22.8 \\
\midrule
\multirow{3}{*}{\textit{City}}
  & 30 & 49.95 & 56.24 & 216.8 & 12.9 \\
  & 42 & 59.26 & 64.47 & 216.3 & 10.8 \\
  & 54 & 64.32 & 70.28 & 210.2 & 12.7 \\
\midrule
\multirow{3}{*}{\textit{Industrial}}
  & 30 & 25.51 & 28.26 & 208.6 & 50.3 \\
  & 42 & 30.19 & 33.11 & 197.0 & 27.5 \\
  & 54 & 33.09 & 36.40 & 194.0 & 53.8 \\
\bottomrule
\end{tabular}

\end{table}

% ---------------------------------------------------------------
\subsection{Cross-Context Comparison}
\label{ssec:crosscase}

\Cref{tab:crosscase} assembles the key performance indicators across
all three study areas to reveal how the trade-off between coverage
retention and runtime reduction varies with urban context.

\begin{table}[t]
  \centering
  \caption{Cross-context performance summary.
           Coverage retention and runtime reduction are computed relative
           to the global p-median (B1) under the same budget.}
  \label{tab:crosscase}
  \begin{tabular}{lrrrrrr}
\toprule
\textbf{Study area}
  & \textbf{Cov.\ (\%)}
  & \textbf{B1 cov.\ (\%)}
  & \textbf{Retention}
  & \textbf{RT (s)}
  & \textbf{B1 RT (s)}
  & \textbf{RT red.} \\
\midrule
Student (Monterrey) & 47.77 & 47.83 & 99.9\% &  25.9 & 78.5 & $+$67.0\% \\
City (Villahermosa) & 59.26 & 60.51 & 97.9\% &  11.2 & 11.4 & $+$1.5\% \\
Industrial (EdoMéx) & 30.19 & 32.38 & 93.2\% &  29.1 & 19.0 & $-$53.1\% \\
\midrule
\multicolumn{7}{l}{\footnotesize Coverage retention = proposed coverage / B1 coverage.} \\
\multicolumn{7}{l}{\footnotesize RT red.\ = (B1 runtime $-$ proposed runtime) / B1 runtime.
                    Negative value means clustering is slower.} \\
\bottomrule
\end{tabular}

\end{table}

The student area is the instance where the proposed framework most clearly
dominates: coverage retention is highest (99.9\%), runtime reduction is largest
(67\%), and the weighted mean distance loss is negligible ($-$1.2~m relative to
the global optimum, meaning the clustered model actually achieves
\emph{shorter} average distances due to the proportional allocation).
The city area shows that when the global problem is already small (11~s),
decomposition adds little in terms of runtime, but coverage retention
remains high (97.9\%) with no meaningful distance cost.
The industrial corridor reveals the limits of fixed-budget decomposition
in supply-constrained, geometrically dispersed settings:
coverage retention drops to 93.2\% and the clustered solve takes longer
than the global one.
The overall message is that territorial decomposition is most effective
in medium-to-large, well-supplied urban instances where candidate-set
density is high enough for subproblems to be genuinely smaller than the
global MILP.
